\documentclass[11pt]{article}
\usepackage{preamble}
\setlength{\parskip}{4pt}

\begin{document}

\daybanner{AP Statistics \textperiodcentered\ Unit 2 \textperiodcentered\ Topic 2.10 \textperiodcentered\ TEACHER KEY}{Unusual Does Not Mean Impossible}

\sectionbanner{CED TETHER}

{\small
\begin{itemize}[leftmargin=1.5em,itemsep=2pt,topsep=2pt]
  \item \textbf{Skill 3.B} --- Using Probability and Simulation: Determine parameters for probability distributions.
  \item \textbf{Skill 4.B} --- Statistical Argumentation: Interpret statistical calculations and findings to assign meaning or assess a claim.
  \item \textbf{EU UNC-3} --- Probabilistic reasoning allows us to anticipate patterns in data.
  \item \textbf{LO UNC-3.C} --- Calculate parameters for a binomial distribution. [Skill 3.B]
  \item \textbf{EK UNC-3.C.1} --- If a random variable is binomial, its mean, $\mu_X$, is np and its standard deviation, $\sigma_X$, is $\surd$(np(1 - p)).
  \item \textbf{LO UNC-3.D} --- Interpret probabilities and parameters for a binomial distribution. [Skill 4.B]
  \item \textbf{EK UNC-3.D.1} --- Probabilities and parameters for a binomial distribution should be interpreted using appropriate units and within the context of a specific population or situation.
\end{itemize}
}
\textbf{Essential question:} How do binomial parameters help distinguish an unusual count from an impossible one?\par
\textbf{DOK-3 skill:} 4.B \textperiodcentered\ \textbf{FRQ pattern:} {\small\texttt{binomial-\allowbreak{}mean-\allowbreak{}sd-\allowbreak{}unusual-\allowbreak{}vs-\allowbreak{}impossible}}\par
\textbf{DOK rationale:} Part (c) weighs competing claims using the day's numerical or graphical evidence and explains a limitation or consequence; the judgment requires linked reasoning beyond a calculation.\par

\textbf{Self-paced bonus sheet.} Students take it when they choose and turn it in when they finish. Everything they need is printed on the sheet. Score part (c) only, E / P / I, by hand --- never AI-graded or auto-scored. (a) and (b) are the ladder up; use them to see where a thin (c) came from.\par

\textbf{Questions to ask}
\begin{itemize}[leftmargin=1.5em,itemsep=2pt,topsep=2pt]
  \item Which number or visible feature supports your claim?
  \item Why does the other claim go beyond that evidence?
  \item What would you tell someone who chose the other claim?
\end{itemize}

\begin{callout}{\IconWarn\ LOOK FOR}{calloutred}
\begin{itemize}[leftmargin=1.5em,itemsep=2pt,topsep=2pt]
  \item A choice with copied numbers but no explanation of how they support it.
  \item Treating a model or average as a guaranteed individual outcome.
\end{itemize}
\end{callout}

\sectionbanner{SCORING PART (c) --- E / P / I}

\textbf{Expected elements}
\begin{itemize}[leftmargin=1.5em,itemsep=2pt,topsep=2pt]
  \item \textbf{reasoning-1} (required) --- Compares 58 with mean 70 and SD about 4.6 in count units
  \item \textbf{reasoning-2} (required) --- Uses the positive 0.0072 chance to distinguish unusual from impossible
  \item \textbf{reasoning-3} (required) --- Supports the analyst and explains the coach's false guaranteed cutoff or proof claim
\end{itemize}

{\small\renewcommand{\arraystretch}{1.25}
\noindent\begin{tabular}{|p{0.5in}|p{5.9in}|}
\hline
\textbf{E} & Includes all three required reasoning elements above, with correct contextual evidence. \\ \hline
\textbf{P} & Includes at least one substantive correct reasoning link above, but omits or misstates another required element. \\ \hline
\textbf{I} & Includes no substantive correct reasoning link above; an unsupported choice or copied number alone is insufficient. \\ \hline
\end{tabular}}

\textbf{Common mistakes}
\begin{itemize}[leftmargin=1.5em,itemsep=2pt,topsep=2pt]
  \item Using 0.70 as the mean count instead of 70
  \item Omitting the square root in the standard deviation
  \item Treating an unusual outcome as impossible or as proof that the model changed
\end{itemize}

\clearpage
\focusbanner{THE PROBLEM --- annotated}

Suppose the same 70 percent free-throw shooter takes 100 attempts under comparable conditions and makes 58. Let $X$ be the number of makes if the career-rate model still applies.\par\smallskip \textbf{Coach:} ``Anything under 60 makes is impossible for a 70 percent shooter, so 58 proves the career rate no longer applies.''\par \textbf{Analyst:} ``A result can be unusual without being impossible; compare 58 with the model's mean and standard deviation.''\par The binomial model gives a chance of about $0.0072$ for 58 or fewer makes; use this provided value.\par
\begin{center}
\textbf{Career-rate model}\par\smallskip
\begin{tabular}{|l|c|c|c|}
\hline
\textbf{Set} & \textbf{Attempts} & \textbf{Rate} & \textbf{Makes} \\ \hline
\textbf{Tonight} & 100 & 0.70 & 58 \\ \hline
\end{tabular}
\end{center}
\begin{firsttakebox}{FIRST TAKE --- one sentence before you work the parts}
Whose statement is more defensible, the coach's or the analyst's? Name one quantity you would check.\par
\answer{Not graded. Listen for whether students use impossible to mean unlikely; the calculation should sharpen that language.}
\end{firsttakebox}

\begin{callout}{\IconBook\ CENTER, SPREAD, AND UNUSUAL COUNTS (keep on the page)}{calloutgreen}
For a binomial count, $\mu_X=np$ and $\sigma_X=\sqrt{np(1-p)}$, both in count units. The mean is an average over many comparable sets of trials. The standard deviation describes a typical distance from that mean. A result far from the mean may be unusual, but a small positive probability does not mean impossible or prove that the model changed.
\end{callout}

\dokbadge{1}~\textbf{(a)} State $n$ and $p$ for this model.\par
\answer{$n=100$ attempts and $p=0.70$ for each attempt.}

\dokbadge{2}~\textbf{(b)} Calculate the binomial mean and standard deviation. Interpret the mean for many sets of 100 shots.\par
\answer{$\mu_X=100(0.70)=70$ makes and $\sigma_X=\sqrt{100(0.70)(0.30)}=\sqrt{21}\approx4.6$ makes. Across many comparable sets of 100 attempts, the model predicts an average of 70 makes.}

\dokbadge{3}~\textbf{(c)} Whose statement is supported? Compare 58 with the mean and standard deviation, and use the given probability to distinguish unusual from impossible. Explain what the coach's claim could make a reader wrongly believe. Write 3--4 sentences.\par
\answer{The analyst is supported: 58 is 12 makes below the mean of 70, more than twice the typical distance of about 4.6 makes. The given chance of 58 or fewer is about $0.0072$, so such a low result is unusual but has positive probability. The coach turns a typical performance level into a guaranteed lower bound. That could make a reader think one unusual set proves the career rate changed, even though chance can produce it.}

\end{document}
